\documentclass{article}
\usepackage[margin=0pt,paperwidth=32cm,paperheight=18cm]{geometry}
\usepackage{tikz}
\pagestyle{empty}

% ==================== 配色 ====================
\definecolor{paper}{RGB}{243,238,224}
\definecolor{skyTop}{RGB}{226,232,230}
\definecolor{blueFar}{RGB}{152,182,192}
\definecolor{blueMid}{RGB}{104,146,158}
\definecolor{greenFar}{RGB}{158,188,156}
\definecolor{greenMid}{RGB}{106,150,108}
\definecolor{greenDeep}{RGB}{56,96,70}
\definecolor{greenDark}{RGB}{26,56,42}
\definecolor{ink}{RGB}{20,30,26}
\definecolor{cinnabar}{RGB}{178,58,46}

% ==================== 元素宏 ====================
\newcommand{\pine}[3]{%
  \begin{scope}[shift={(#1,#2)}, scale=#3]
    \draw[ink, line width=0.5pt]
      (0,0) .. controls (0.08,0.5) and (-0.06,1.0) .. (0.02,1.55);
    \foreach \y/\w in {1.45/0.28, 1.20/0.40, 0.95/0.50, 0.68/0.46, 0.42/0.38, 0.18/0.28}{
      \draw[ink!70, line width=0.25pt] (0.01,\y) -- (\w,\y-0.14);
      \draw[ink!70, line width=0.25pt] (0.01,\y) -- (-\w,\y-0.14);
      \fill[greenDark, opacity=0.85] (\w,\y-0.18) circle (0.055);
      \fill[greenDark, opacity=0.85] (-\w,\y-0.18) circle (0.055);
    }
    \fill[greenDark, opacity=0.9] (0.02,1.60) circle (0.07);
  \end{scope}
}

\newcommand{\broadleaf}[3]{%
  \begin{scope}[shift={(#1,#2)}, scale=#3]
    \draw[ink!75, line width=0.4pt]
      (0,0) .. controls (0.08,0.25) and (-0.05,0.5) .. (0.03,0.70);
    \foreach \dx/\dy in {-0.28/0.72, -0.15/0.88, 0.02/0.94, 0.18/0.86, 0.30/0.70,
                          0.12/0.66, -0.20/0.60, -0.02/0.76, 0.22/0.52, -0.30/0.45}{
      \fill[greenDark, opacity=0.66] (\dx,\dy) circle (0.13);
    }
  \end{scope}
}

\newcommand{\willow}[3]{%
  \begin{scope}[shift={(#1,#2)}, scale=#3]
    \draw[ink!75, line width=0.4pt]
      (0,0) .. controls (0.05,0.3) and (-0.02,0.55) .. (0.05,0.80);
    \foreach \x in {-0.45,-0.30,-0.15,0,0.15,0.30,0.45}{
      \draw[greenMid!70, opacity=0.55, line width=0.25pt]
        (0.05,0.82) .. controls (\x,0.78) and (\x*0.85,0.55) .. (\x*0.70,0.28);
    }
    \fill[greenDark, opacity=0.65] (0.05,0.84) circle (0.15);
  \end{scope}
}

\newcommand{\pavilion}[3]{%
  \begin{scope}[shift={(#1,#2)}, scale=#3]
    \fill[ink, opacity=0.88]
      (-0.70,0.52) .. controls (-0.45,0.86) and (0.45,0.86) .. (0.70,0.52)
      .. controls (0.42,0.60) and (-0.42,0.60) .. cycle;
    \draw[ink, line width=0.5pt]
      (-0.70,0.52) .. controls (-0.45,0.86) and (0.45,0.86) .. (0.70,0.52);
    \fill[ink, opacity=0.8] (-0.36,0.12) rectangle (-0.29,0.52);
    \fill[ink, opacity=0.8] ( 0.29,0.12) rectangle ( 0.36,0.52);
    \fill[ink, opacity=0.8] (-0.29,0.12) rectangle (-0.12,0.19);
    \fill[ink, opacity=0.8] ( 0.12,0.12) rectangle ( 0.29,0.19);
    \fill[ink] (0,0.84) -- (0.05,0.94) -- (-0.05,0.94) -- cycle;
    \draw[ink, line width=0.35pt] (0,0.94) -- (0,1.06);
  \end{scope}
}

\newcommand{\boat}[3]{%
  \begin{scope}[shift={(#1,#2)}, scale=#3]
    \fill[ink, opacity=0.72]
      (-0.45,0) .. controls (-0.20,-0.14) and (0.20,-0.14) .. (0.45,0)
      .. controls (0.20,0.06) and (-0.20,0.06) .. cycle;
    \draw[ink, opacity=0.6, line width=0.3pt] (0.02,0.05) -- (0.02,0.62);
    \fill[ink, opacity=0.32]
      (0.04,0.58) .. controls (0.28,0.40) and (0.28,0.18) .. (0.04,0.12) -- cycle;
  \end{scope}
}

% ==================== 山峦轮廓坐标 ====================
\def\peakA{(0,10.0) (2,11.5) (4,12.8) (6,11.6) (8,13.0) (10,13.6) (12,12.4) (14,13.4) (16,12.2) (18,13.0) (20,11.8) (22,12.4) (24,11.0) (26,11.4) (28,10.2) (30,10.6) (32,9.8)}
\def\peakB{(0,8.8) (2,10.2) (4,11.6) (6,10.4) (8,11.8) (10,12.4) (12,11.2) (14,12.2) (16,11.0) (18,11.8) (20,10.6) (22,11.2) (24,9.8) (26,10.2) (28,9.0) (30,9.4) (32,8.6)}
\def\peakC{(0,7.5) (1.5,9.0) (3,10.4) (4.5,9.2) (6,10.6) (7.5,11.2) (9,10.0) (10.5,10.8) (12,9.6) (13.5,10.4) (15,9.2) (16.5,9.8) (18,8.6) (19.5,9.0) (21,7.8) (23,8.2) (25,7.0) (27,7.4) (29,6.4) (32,6.8)}
\def\peakD{(0,5.8) (1.5,7.2) (3,8.6) (4.5,7.4) (6,8.8) (7.5,9.4) (9,8.2) (10.5,8.9) (12,7.8) (13.5,8.5) (15,7.3) (16.5,7.9) (18,6.7) (19.5,7.2) (21,6.0) (23,6.5) (25,5.4) (27,5.8) (29,4.8) (32,5.2)}
\def\peakE{(0,4.0) (1.5,5.5) (3,7.2) (4.5,6.0) (6,7.4) (7.5,8.0) (9,6.8) (10.5,7.4) (12,6.2) (13.5,6.9) (15,5.7) (16.5,6.2) (18,5.0) (19.5,5.4) (21,4.4) (23,4.7) (25,3.8) (27,4.1) (29,3.4) (32,3.7)}
\def\peakF{(0,2.8) (1.5,3.8) (3,4.9) (4.5,3.9) (6,4.9) (7.5,5.3) (9,4.4) (10.5,4.8) (12,3.9) (14,4.2) (16,3.4) (18,3.6) (20,3.0) (22,3.2) (24,2.9) (32,2.6)}

\begin{document}
\noindent
\begin{tikzpicture}[remember picture, overlay]
  \path (current page.south west) coordinate (SW);
  \begin{scope}[shift={(SW)}, x=1cm, y=1cm]

    % ==================== 天空 + 纸底 ====================
    \shade[top color=skyTop, bottom color=paper] (0,0) rectangle (32,18);

    % 纸纹
    \begin{scope}
      \clip (0,0) rectangle (32,18);
      \foreach \x in {0,0.75,...,32}
        \foreach \y in {0,0.75,...,18}
          \fill[ink, opacity=0.020] (\x,\y) circle (0.018);
    \end{scope}

    % ==================== 远山 A（淡石青） ====================
    \fill[blueFar, opacity=0.28]
      plot[smooth, tension=0.55] coordinates {\peakA}
      -- (32,2.6) -- (0,2.6) -- cycle;

    % ==================== 远山 B ====================
    \fill[blueMid, opacity=0.42]
      plot[smooth, tension=0.55] coordinates {\peakB}
      -- (32,2.6) -- (0,2.6) -- cycle;

    % ---------- 云雾带 1 ----------
    \fill[paper, opacity=0.55] (0,8.4) rectangle (32,9.1);
    \fill[paper, opacity=0.30] (0,9.4) rectangle (32,9.8);

    % ==================== 中山 C（石绿） ====================
    \fill[greenFar, opacity=0.60]
      plot[smooth, tension=0.55] coordinates {\peakC}
      -- (32,2.6) -- (0,2.6) -- cycle;

    % 中山 C 山脊
    \foreach \x/\y in {3/10.4, 6/10.6, 7.5/11.2, 10.5/10.8, 13.5/10.4, 16.5/9.8}{
      \draw[greenDark, opacity=0.22, line width=0.4pt]
        (\x,\y) .. controls (\x-0.15,\y-0.6) and (\x+0.20,\y-1.1) .. (\x-0.02,\y-1.8);
    }

    % ==================== 中山 D ====================
    \fill[greenMid, opacity=0.78]
      plot[smooth, tension=0.55] coordinates {\peakD}
      -- (32,2.6) -- (0,2.6) -- cycle;

    % 中山 D 山脊
    \foreach \x/\y in {3/8.6, 6/8.8, 7.5/9.4, 10.5/8.9, 13.5/8.5, 16.5/7.9, 19.5/7.2}{
      \draw[greenDark, opacity=0.28, line width=0.45pt]
        (\x,\y) .. controls (\x-0.18,\y-0.5) and (\x+0.22,\y-1.0) .. (\x,\y-1.7);
    }

    % 中山 D 皴法
    \begin{scope}
      \clip plot[smooth, tension=0.55] coordinates {\peakD}
            -- (32,2.6) -- (0,2.6) -- cycle;
      \foreach \x/\y in {
        1.5/7.0, 2.4/7.6, 3.3/7.4, 4.2/6.8, 5.2/7.8, 6.1/7.6,
        7.0/8.4, 8.0/7.4, 8.9/7.6, 9.8/7.8, 10.8/7.6, 11.7/6.8,
        12.6/7.4, 13.5/7.4, 14.5/6.4, 15.5/6.9, 16.5/6.7, 17.5/5.9,
        18.5/6.1, 19.5/6.0, 20.5/5.1, 22/5.4, 24/4.6
      }{
        \draw[greenDark, opacity=0.20, line width=0.3pt]
          (\x,\y) .. controls (\x+0.14,\y-0.35) and (\x-0.08,\y-0.65) .. (\x+0.06,\y-1.0);
      }
    \end{scope}

    % ---------- 云雾带 2 ----------
    \fill[paper, opacity=0.62] (0,5.4) rectangle (32,6.0);
    \fill[paper, opacity=0.40] (0,6.2) rectangle (32,6.6);
    \fill[paper, opacity=0.20] (0,6.9) rectangle (32,7.2);

    % ==================== 近山 E ====================
    \fill[greenDeep, opacity=0.90]
      plot[smooth, tension=0.55] coordinates {\peakE}
      -- (32,2.6) -- (0,2.6) -- cycle;

    % 近山 E 山脊
    \foreach \x/\y in {3/7.2, 6/7.4, 7.5/8.0, 10.5/7.4, 13.5/6.9, 16.5/6.2, 19.5/5.4}{
      \draw[greenDark, opacity=0.42, line width=0.55pt]
        (\x,\y) .. controls (\x-0.20,\y-0.45) and (\x+0.25,\y-0.95) .. (\x,\y-1.5);
    }

    % 近山 E 皴法
    \begin{scope}
      \clip plot[smooth, tension=0.55] coordinates {\peakE}
            -- (32,2.6) -- (0,2.6) -- cycle;
      \foreach \x/\y in {
        1.2/4.8, 2.0/5.6, 2.8/6.2, 3.6/5.8, 4.4/5.2, 5.2/6.2, 6.0/6.4,
        6.8/6.8, 7.6/6.4, 8.4/5.9, 9.2/6.1, 10.0/6.3, 10.8/6.0, 11.6/5.3,
        12.4/5.7, 13.2/5.6, 14.0/4.9, 14.8/5.2, 15.6/4.7, 16.4/4.7,
        17.2/4.1, 18.0/4.2, 18.8/3.7, 19.6/3.8, 20.4/3.3
      }{
        \draw[greenDark, opacity=0.28, line width=0.3pt]
          (\x,\y) .. controls (\x+0.15,\y-0.30) and (\x-0.10,\y-0.60) .. (\x+0.05,\y-0.95);
      }
    \end{scope}

    % ==================== 近岸 F（墨绿） ====================
    \fill[greenDark, opacity=0.94]
      plot[smooth, tension=0.55] coordinates {\peakF}
      -- (32,2.6) -- (0,2.6) -- cycle;

    % 苔点
    \foreach \x/\y in {
      0.8/3.4, 1.5/4.2, 2.3/3.8, 3.1/4.6, 3.9/3.5,
      4.6/4.4, 5.3/5.0, 6.1/4.2, 6.9/5.2, 7.6/4.6,
      8.3/5.4, 9.1/4.4, 9.8/5.2, 10.6/4.6, 11.3/5.4,
      12.1/4.3, 12.8/5.0, 13.6/4.2, 14.4/4.8, 15.1/3.8,
      16.0/4.4, 16.8/3.6, 17.6/4.0, 18.4/3.2, 19.2/3.6,
      20.0/3.0, 20.8/3.4, 21.6/2.9
    }{
      \fill[greenDark!60!black, opacity=0.30] (\x,\y) circle (0.04);
    }

    % ==================== 瀑布 ====================
    \begin{scope}
      \clip plot[smooth, tension=0.55] coordinates {\peakD}
            -- (32,2.6) -- (0,2.6) -- cycle;
      \fill[paper, opacity=0.80]
        (11.55,9.0) .. controls (11.75,8.2) and (11.45,7.2) .. (11.65,6.2)
        .. controls (11.85,5.2) and (11.55,4.4) .. (11.75,3.6)
        -- (12.35,3.6) .. controls (12.25,4.4) and (12.55,5.2) .. (12.35,6.2)
        .. controls (12.15,7.2) and (12.45,8.2) .. (12.25,9.0) -- cycle;
      \foreach \y in {3.8,4.4,5.0,5.6,6.2,6.8,7.4,8.0,8.6}{
        \draw[paper, opacity=0.95, line width=0.4pt]
          (11.70,\y) .. controls (11.95,\y+0.06) and (12.05,\y-0.06) .. (12.30,\y);
      }
    \end{scope}

    % ==================== 水面 ====================
    \fill[blueFar, opacity=0.20] (0,0) rectangle (32,2.6);
    \fill[paper, opacity=0.35] (0,2.4) rectangle (32,2.75);

    % 倒影
    \foreach \y/\op in {0.30/0.28, 0.70/0.24, 1.10/0.20, 1.50/0.16, 1.90/0.12, 2.25/0.09}{
      \draw[greenDeep, opacity=\op, line width=0.6pt]
        (0.5,\y) .. controls (5,\y+0.08) and (10,\y-0.08) .. (16,\y)
                 .. controls (20,\y+0.06) and (24,\y-0.06) .. (28,\y);
    }

    % 水纹
    \foreach \x/\y/\w in {3/1.2/2.0, 7/0.6/2.5, 12/1.8/1.8, 16/0.9/2.2, 21/1.5/1.6, 25/0.5/2.0}{
      \draw[blueMid, opacity=0.28, line width=0.4pt] (\x,\y) -- (\x+\w,\y);
    }

    % ==================== 松林（左下） ====================
    \pine{1.00}{3.05}{1.10}
    \pine{2.20}{3.45}{0.95}
    \pine{3.40}{3.85}{0.82}
    \pine{4.40}{4.20}{0.70}
    \pine{5.30}{4.55}{0.58}

    % ==================== 杂树 ====================
    \broadleaf{6.6}{4.4}{0.90}
    \broadleaf{7.7}{4.9}{0.75}
    \broadleaf{8.7}{5.2}{0.62}

    % ==================== 垂柳（水岸） ====================
    \willow{13.6}{3.3}{0.95}
    \willow{14.6}{3.6}{0.80}

    % ==================== 亭子 ====================
    \pavilion{16.5}{3.2}{1.00}

    % ==================== 小船 ====================
    \boat{22.0}{1.4}{1.00}
    \boat{25.5}{0.9}{0.70}

    % ==================== 飞鸟 ====================
    \foreach \x/\y/\s in {
      18.5/15.5/0.22, 19.7/15.8/0.18, 20.8/15.4/0.16,
      21.8/15.7/0.14, 22.8/15.3/0.12, 23.6/15.6/0.10
    }{
      \draw[ink, opacity=0.38, line width=0.55pt]
        (\x,\y) .. controls (\x+0.12,\y+\s) and (\x+0.22,\y+\s) .. (\x+0.32,\y)
                .. controls (\x+0.42,\y+\s) and (\x+0.52,\y+\s) .. (\x+0.64,\y);
    }
    \foreach \x/\y/\s in {14.5/16.4/0.18, 15.5/16.7/0.15, 16.4/16.3/0.13}{
      \draw[ink, opacity=0.28, line width=0.45pt]
        (\x,\y) .. controls (\x+0.10,\y+\s) and (\x+0.18,\y+\s) .. (\x+0.26,\y)
                .. controls (\x+0.34,\y+\s) and (\x+0.42,\y+\s) .. (\x+0.52,\y);
    }

    % ==================== 朱砂印章 ====================
    \fill[cinnabar, opacity=0.88] (29.20,1.10) rectangle (30.90,2.80);
    \draw[paper, opacity=0.90, line width=0.6pt] (29.20,1.95) -- (30.90,1.95);
    \draw[paper, opacity=0.90, line width=0.6pt] (30.05,1.10) -- (30.05,2.80);
    \draw[paper, opacity=0.45, line width=0.3pt] (29.38,1.28) rectangle (30.72,2.62);

    % ==================== 竖排题款 ====================
    \foreach \i/\ch in {
      0/Q, 1/I, 2/N, 3/G, 4/L, 5/U, 6/S, 7/H, 8/A, 9/N,
      10/S, 11/H, 12/U, 13/I
    }{
      \node[anchor=center, text=ink, opacity=0.45,
            font=\fontsize{7}{7}\selectfont]
        at (29.95, 7.45 - \i*0.34) {\ch};
    }

  \end{scope}
\end{tikzpicture}
\end{document}
